\documentclass[10pt,letterpaper]{article}

\usepackage[letterpaper,margin=0.72in]{geometry}
\usepackage[T1]{fontenc}
\usepackage[utf8]{inputenc}
\usepackage{lmodern}
\usepackage{microtype}
\usepackage{enumitem}
\usepackage{xcolor}
\usepackage{hyperref}
\usepackage{titlesec}
\usepackage{tabularx}
\usepackage[backend=biber,style=numeric,sorting=ydnt,maxbibnames=99,giveninits=true,defernumbers=true]{biblatex}

\addbibresource{../rob_papers.bib}

\definecolor{ink}{HTML}{24211D}
\definecolor{muted}{HTML}{6B6259}
\definecolor{accent}{HTML}{176B7A}
\definecolor{rule}{HTML}{D8D0C5}
\definecolor{soft}{HTML}{F5F1EA}

\hypersetup{
  colorlinks=true,
  urlcolor=accent,
  linkcolor=accent,
  citecolor=accent,
  pdfauthor={Robert M. Gower},
  pdftitle={Robert M. Gower -- Curriculum Vitae}
}

\pagestyle{empty}
\setlength{\parindent}{0pt}
\setlist[itemize]{leftmargin=1.15em,itemsep=0.18em,topsep=0.18em}
\setlist[enumerate]{leftmargin=1.4em,itemsep=0.18em,topsep=0.18em}
\renewcommand{\arraystretch}{1.08}

\titleformat{\section}
  {\Large\bfseries\color{ink}}
  {}
  {0pt}
  {}
  [\vspace{0.25em}{\color{rule}\titlerule}]
\titlespacing*{\section}{0pt}{1.25em}{0.55em}

\newcommand{\cvtag}[1]{%
  \begingroup
  \setlength{\fboxsep}{3pt}%
  \colorbox{soft}{\textcolor{accent}{\footnotesize\bfseries #1}}%
  \endgroup
}

\newcommand{\entry}[4]{%
  \begin{tabularx}{\textwidth}{@{}p{0.19\textwidth}X@{}}
    \textcolor{muted}{#1} & \textbf{#2} #3\\
    & {\small\textcolor{muted}{#4}}
  \end{tabularx}\vspace{0.35em}
}

\newcommand{\simpleentry}[2]{%
  \begin{tabularx}{\textwidth}{@{}p{0.19\textwidth}X@{}}
    \textcolor{muted}{#1} & #2
  \end{tabularx}\vspace{0.28em}
}

\newcommand{\talkentry}[4]{%
  \simpleentry{#1}{\textbf{#2}. #3\if\relax\detokenize{#4}\relax\else\hfill \href{#4}{recording/link}\fi}
}

% Compatibility for the older included CV fragments.
\newenvironment{bibsection}
  {\begin{description}[leftmargin=0.19\textwidth,style=nextline,itemsep=0.28em,topsep=0.2em]}
  {\end{description}}
\newenvironment{innerlist}
  {\begin{itemize}[leftmargin=1.15em,itemsep=0.12em,topsep=0.12em]}
  {\end{itemize}}

\begin{document}

{\Huge\bfseries Robert M. Gower}\\[-0.15em]
{\Large\color{muted}Research Scientist, Flatiron Institute}\\[0.65em]
{\color{rule}\hrule}
\vspace{0.7em}

\begin{tabularx}{\textwidth}{@{}Xr@{}}
\href{mailto:gowerrobert@gmail.com}{gowerrobert@gmail.com} \quad
\href{https://gowerrobert.github.io}{gowerrobert.github.io} \quad
\href{https://scholar.google.com/citations?user=okKw87MAAAAJ\&hl=en}{Google Scholar} \quad
\href{https://github.com/gowerrobert}{GitHub} \quad
\href{https://www.linkedin.com/in/robert-mansel-gower-a27bba23}{LinkedIn}
&
Updated \today
\end{tabularx}

\vspace{0.8em}
\cvtag{modern optimization}
\cvtag{adaptive methods}
\cvtag{stochastic optimization}
\cvtag{machine learning}
\cvtag{randomized numerical linear algebra}

\section{Profile}
I work on the mathematics of modern optimization: why methods such as Adam, Muon, Polyak step-sizes, and non-Euclidean gradient descent work, fail, and sometimes come up with new methods. My current research sits between optimization theory and machine-learning practice, with a focus on simple algorithms that exploit structure. Previously, I worked on randomized numerical linear algebra, automatic differentiation, and quasi-Newton methods.

\section{Appointments and Education}
\entry{11/2021--present}{Research Scientist}{, Center for Computational Mathematics, Flatiron Institute}{New York, NY}
\entry{06/2021--11/2021}{Visiting Scientist}{, Google Research}{New York, NY}
\entry{09/2017--06/2021}{Assistant Professor in Machine Learning}{, Telecom Paris, Institut Polytechnique de Paris}{Paris, France}
\entry{01/2020--08/2020}{Visiting Scientist}{, Facebook Research}{New York, NY}
\entry{08/2016--09/2017}{Laureate Fellow}{, Fondation Sciences Mathematiques de Paris, ENS and Inria}{Paris, France}
\entry{09/2012--06/2016}{Ph.D.}{, Optimization and Operations Research, University of Edinburgh}{Thesis: \emph{Sketch and Project: Randomized Iterative Methods for Linear Systems and Inverting Matrices}}
\entry{03/2009--04/2011}{M.Sc.}{, Applied Mathematics, University of Campinas}{Thesis: \emph{Hessian Matrices via Automatic Differentiation}}
\entry{03/2005--03/2009}{B.Sc.}{, Applied and Computational Mathematics, University of Campinas}{Brazil}

\section{Selected Recent News}
\simpleentry{2026}{\href{https://blog.iclr.cc/2026/04/23/announcing-the-iclr-2026-outstanding-papers/}{\textbf{The Polar Express}} received the ICLR 2026 Outstanding Paper Honorable Mention and was presented as an oral.}
\simpleentry{2026}{Upcoming talk at the \href{https://www.siam.org/conferences-events/siam-conferences/op26/}{SIAM Conference on Optimization (OP26)}, June 2--5, 2026.}
\simpleentry{2025}{\href{https://neurips.cc/virtual/2025/loc/mexico-city/oral/130313}{\textbf{In Search of Adam's Secret Sauce}} was an oral presentation at NeurIPS 2025.}

\section{Prizes, Honors, and Funding}
\begin{itemize}
\item ICLR 2026 Outstanding Paper Honorable Mention for \emph{The Polar Express: Optimal Matrix Sign Methods and Their Application to the Muon Algorithm}.
\item NeurIPS 2025 Top Reviewer.
\item ICML 2021 Best Reviewer, Top 10\%.
\item NeurIPS 2020 Top 10\% reviewer.
\item Thomas Jefferson Fund grant, ``Matrix equations for big data compression'', \$20k, 2019.
\item CIFRE Facebook Research funding, \$72k, 2019.
\item ICML 2019 Top 5\% reviewer.
\item LabEx Mathematique Hadamard funding, ``Scalable stochastic variance reduced gradient methods'', 8k EUR, 2018.
\item Leslie Fox Prize for Numerical Analysis, 2nd place, 2017.
\item Fondation Sciences Mathematiques de Paris Fellowship at Inria/ENS, 122k EUR, 2016.
\item Best Talk Prize, Irish SIAM student meeting, 2014.
\item University of Edinburgh Ph.D. scholarship, 85.1k GBP, 2012.
\end{itemize}

\section{People and Mentoring}
\begin{itemize}
\item \href{https://ngazagna.github.io/}{Nidham Gazagnadou}, Ph.D. student, 2018--2021; now Research Scientist at Sony AI.
\item \href{https://rui-yuan91.github.io/}{Rui Yuan}, Ph.D. student, 2019--2023; now AI Research Scientist at Stellantis.
\item \href{https://sjelassi.github.io/}{Samy Jelassi}, intern, 2017; Ph.D. at Princeton.
\item \href{https://www.cs.cornell.edu/~siyimeng/}{Si Yi Meng (Cathy)}, intern, 2022 and 2023.
\item \href{https://fabian-sp.github.io/}{Fabian Schaipp}, guest researcher, 2022 and 2023.
\item \href{https://cispa.de/en/people/c01slha}{Slavomir Hanzely}, intern, 2022; now Researcher at CISPA.
\item \href{https://cs.stanford.edu/~amishkin/}{Aaron Mishkin}, intern, 2023; now Postdoc at EPFL.
\item \href{https://davpersson.github.io/}{David Persson}, Flatiron Fellow, 2024--present.
\item \href{https://mtcrawshaw.github.io/}{Michael Crawshaw}, intern, 2025; joining Flatiron as a Fellow in 2026.
\item \href{https://parshakova.github.io/}{Tetiana Parshakova}, Flatiron Fellow, 2025--present.
\end{itemize}

\section{Recorded Talks}
\talkentry{2026}{The Polar Express}{Instituto de Matematica Pura e Aplicada}{https://www.youtube.com/live/phM8GzT7ZV4}
\talkentry{2024}{New viewpoints, variants and convergence theory for stochastic Polyak step-sizes}{MTL MLOpt}{https://www.youtube.com/watch?v=7lY-SQKslwk}
\talkentry{2019}{SGD: General Analysis and Improved Rates}{ICML oral presentation}{https://icml.cc/virtual/2019/oral/4669}

\section{Selected Invited Talks}
	\begin{bibsection}
	\item [Over {\bf 30 talks}] in international conferences, seminars and workshops including.
\item [2019] Parietal team seminar, Inria, Saclay
\item  International Conference on Continuous Optimization (ICCOPT), invited session, Berlin.
\item Journ\'ee Statistique et Apprentissage, L'Institut des Hautes \'Etudes Scientifiques.
\item [2018] The 21st International Conference on
Artificial Intelligence and Statistics (AISTATS), {\bf semi plenary presentation} and two posters, Lanzarote, Canary Islands


\item {Séminaire Parisien d'Optimisation},  Institut Henri Poincaré

\item {International Symposium of Mathematical Optimization}, invited session on Fast Converging Stochastic Optimization Algorithms, Bordeaux


\item Seminar for the Machine Learning and Optimization group, {EPFL} %,  26th of June, 2018
		
		\item [2017] \href{https://www.fondation-hadamard.fr/fr/pgmo/pgmodays}{PGMO days}, EDF'Lab Paris-Saclay
		\item Thirty-first Conference on Neural Information Processing Systems (Neurips), Poster, Vancouver.
		
		\item \href{http://optimization2017.fc.ul.pt/}{Optimization 2017}, Faculdade de Ciências of the Universidade de Lisboa, invited talk, Lisboa

\item \href{https://numericalanalysisconference.org.uk/meetings/2017}{The 27th Biennial Numerical Analysis Conference}, Randomized Iterative Methods for Computing the Pseudoinverse, University of Strathclyde, Glasgow

\item \href{http://people.maths.ox.ac.uk/wathen/fox/prize.php}{18th IMA Leslie Fox Prize talk in Numerical Analysis}, awarded 2nd place, University of Strathclyde, Glasgow

%\item SIAM Optimization 2017, Vancouver, Canada

\item  \href{http://fjml.marcocuturi.net/speakers/}{France/Japan Machine Learning Workshop}, Ec\'ole Normale Sup\'erieure, Paris
		\item	\href{https://thoth.inrialpes.fr/people/mairal/macaron/workshop2017/}{Optimization, machine learning, and pluri-disciplinarity workshop}, Paris, Inria Grenoble Rhone-Alpes, September 21-22.
		\item Probabilistic Numerics seminar, Max-Planck-Institute Tubingen
		\item Operations Research seminar, Center for Operations Research and Econometrics, Universit\'e catholique de Louvain
		\item [2016]		
		 Computational Mathematics and Applications seminar, Rutherford Appleton Laboratory
		\item T\'el\'ecom ParisTech, Machine Learning and Statistics seminar, Paris
 	    \item International Conference on Continuous Optimization, Tokyo
 	     \item  5th IMA Conference on Numerical Linear Algebra and Optimization, University of Birmingham
        International Conference on Machine Learning ({\bf ICML}), New York
		\item Cambridge Image Analysis group seminar, University of Cambridge \\
		 \emph{Randomized iterative methods for linear systems}
\item  Inria, SIERRA team seminar, Paris,
\emph{Randomized iterative methods for linear systems and inverting matrices}
		\item[ 2015]  	\href{http://www.icms.org.uk/workshop.php?id=367}{Distributed machine learning and optimization}, Alan Turing Scoping Workshop, Edinburgh, UK \\
		\emph{ Distributed Randomized Iterative Methods for Linear Systems }
\item The International Symposium on Mathematical Programming, Pittsburgh, USA. \\
 \emph{Action constrained quasi-Newton methods}
\item Optimization and Big Data 2015, Edinburgh, UK\\ 
 \emph{Randomized iterative methods for linear systems and inverting matrices}
\item[ 2014] 
 \href{http://www.maths.nuigalway.ie/SIAM-Galway/PG-Conf-14/}{Irish Applied Mathematics Research Students' Meeting}, Galway, Ireland.\\
 \emph{Hunting inverses of matrix fields}  \\
  {\bf [Best Talk Prize]} 
 \item \href{https://info.maths.ed.ac.uk/pg-pages/colloquium/colloquium_1415_1.html}{Postgraduate Mathematics Colloquium}, Edinburgh, UK\\
\emph{The history of optimization in blood and booze}
\item	  Cambridge Mathematics Society Research in the UK Afternoon, Cambridge, UK \\  \emph{Unconstrained optimization methods}
 \item \href{http://europt2014.univ-perp.fr/}{European Workshop
on Advances in Continuous Optimization, Perpignan}. \\ 
 \emph{Generalizations of the quasi-Newton methods using an image constraint}
\item[ 2013]
 \href{http://www.norg.uminho.pt/iccopt2013-SP-contents/wed_b_23.htm}{ International Conference on Continuous Optimization}, Lisbon, Portugal \\
\emph{Third order methods and third order derivatives} 
\item EURO-INFORMS Joint International Meeting, Rome, Italy \\
 \emph{Third order methods using slices of the tensor and AD developments}
\item[ 2011]
28th Brazilian Colloquium on Mathematics, IMPA, Rio de Janeiro, Brazil \\ \emph{Automatic differentiation of Hessian matrices}
\end{bibsection}	
% 
%
%\section{Poster Sessions}
%\begin{bibsection}
%\item[ 2015] Data Science Research Day, Edinburgh, UK\\
%\emph{Randomized methods for inverting matrices}
%\item[ 2014] \href{http://www.maths.ed.ac.uk/~rtappend/CLAODE_Workshop/home.html}{Computational Linear Algebra and Optimization for the Digital Economy}, Edinburgh\\
%\emph{Automatic high order tensor derivatives}
%\item[ 2011]
%Brazilian Society on Applied and Computational Mathematics\\
% \'Aguas de Lind\'oia, Brasil\\
%\emph{A new automatic differentiation mode for sparse Hessian matrices}
%\end{bibsection}	

\section{Teaching}
I have given over 360 hours of lectures on optimization, numerical analysis, and machine learning. Prior to that, as a graduate student at the University of Edinburgh, I gave over 350 hours of tutorials and lectures.
\begin{itemize}
\item 2019--2020: \emph{Machine Learning in High Dimensions}, graduate course, Telecom Paris.
\item 2017--2020: \emph{Optimization and Numerical Analysis}, undergraduate course, Telecom Paris.
\item 2017--2020: \emph{Optimization for Data Science}, Data Science Master's, Institut Polytechnique de Paris.
\item 2019: \emph{Stochastic Optimization for Machine Learning}, African Master's of Machine Intelligence at AIMS, Rwanda.
\item 2017--2019: Summer school lectures on stochastic optimization for machine learning at Yerevan State University and the University of Novi Sad.
\item 2010--2015: Over 350 hours of tutorials and lectures at the University of Edinburgh, including linear algebra, calculus, financial mathematics, operations research, optimization theory, proofs, and problem solving.
\end{itemize}

\section{Service}
Editorial work and reviewing for:
\begin{innerlist}
\item[] Action Editor, \emph{Transactions of Machine Learning Research} (TMLR)
\item[] Editorial Board Reviewers, \emph{Journal of Machine Learning Research} (JMLR)
\item[] Area Chair / Meta-Reviewer, \emph{International Conference on Machine Learning} (ICML), 2025
\item[] Area Chair, \emph{Conference on Neural Information Processing Systems} (NeurIPS), 2021 and 2022
\item[] Senior Program Committee, \emph{International Joint Conference on Artificial Intelligence} (IJCAI), 2020
\item[] Reviewer, \emph{Conference on Neural Information Processing Systems} (NeurIPS)
\item[] \emph{International Conference on Machine Learning} (ICML)
\item[] \emph{Journal of Machine Learning Research} (JMLR)
\item[] \emph{Proceedings of the IEEE}
\item[] \emph{IEEE Transactions on Signal Processing}
\item[] \emph{SIAM Journal on Scientific Computing} (SISC)
\item[] \emph{SIAM Journal on Optimization} (SIOPT)
			\item[] \emph{Optimization Methods and Software}, Taylor \& Francis
			\item[] \emph{Computational and Applied Mathematics}, Springer
			\item[] \emph{Mathematical Programming}, Springer Computation
			\item[] \emph{BIT Numerical Mathematics}, Springer
			\item[] \emph{Numerical Algorithms}, Springer
		\end{innerlist}	


\section{Publications}
All entries below are generated directly from \texttt{../rob\_papers.bib}.
\nocite{*}
\printbibliography[heading=none]

\end{document}
